\documentclass[10pt,a4paper]{amsart}
\usepackage[margin=19mm]{geometry}
\usepackage{amsmath,amssymb,mathtools}
\usepackage[scr=rsfs,cal=euler]{mathalfa}
\usepackage{xfrac}
\usepackage[unicode,hidelinks,bookmarks=false]{hyperref}
\newcommand{\ie}{i.e.\ }
\newcommand{\cf}{cf.\ }
\newcommand{\resp}{respectively\ }
\newcommand{\Subsection}{\subsection}
\providecommand{\nobreakdash}{\nobreak\hbox{-}}
\setlength{\emergencystretch}{2em}
\multlinegap0pt
\usepackage{amssymb}
\usepackage{enumerate}
\usepackage{caption}
\usepackage[shortlabels]{enumitem}
\usepackage{booktabs}
\usepackage{algorithm}\usepackage{algpseudocode}
\usepackage{float}
\usepackage[all]{xy}
\usepackage{tikz}
\newtheorem{thm}{Theorem}
\usepackage{cleveref}
\crefname{thm}{Theorem}{Theorems}
\DeclareMathOperator{\Sat}{Sat}
\newcommand{\Z}{{\mathbb Z}}
\title{Upper Bounds for Sequence Saturation}
\author{Shihan Kanungo}
\date{Source: arXiv:2512.17683v1 (CC BY 4.0)}
\begin{document}
\maketitle

\noindent\emph{Source and licence.} Upper Bounds for Sequence Saturation. Version arXiv:2512.17683v1 (CC BY 4.0), distributed by arXiv under the Creative Commons Attribution 4.0 International licence, http://creativecommons.org/licenses/by/4.0/, as recorded in the arXiv licence metadata for this exact version and shown on the abstract page https://arxiv.org/abs/2512.17683v1. Full licence notice: \texttt{licenses/arxiv-2512.17683-CC-BY-4.0.txt}.\par\smallskip

\noindent\textit{Editorial scope.} Counted unit: the article's thm sequences (source0.tex lines 737--739). The article's complete original proof of it is delivered with it (source0.tex lines 740--760, one \texttt{proof} environment of 21 lines). No mathematics is written, repaired or re-derived: the statement and the proof are the article's own lines, in the article's own notation, and no retained line is substituted or re-flowed.  Retained uncounted as the source-local context the proof needs: lines 691--693.  Nothing was omitted from any retained span, so the package carries no omission record.  No retained line contains a citation, so the wrapper carries no bibliography and no bibliography entry.  No retained line contains a figure environment, an image-inclusion command or drawing code, so no image is delivered and the package declares no asset dependency.  Editorial formatting only, in addition: an AMS \texttt{amsart} preamble that restates the article's own declarations and macros used by the retained lines, namely the environments \texttt{thm} on separate counters, together with the standard packages those lines need.  The retained environments print in this document as Theorem 1 in order of appearance, whereas the article prints the same items with the numbering of its own full document.  The complete native source of the article as distributed in the arXiv e-print for this exact version is bundled unchanged as \texttt{sources/191337/source0.tex}.
\begin{thm}\label{thm: aa...bb}
If $u$ is irreducible and of the form $aa\dots bb$, then $\Sat(n,u)=O(n)$.
\end{thm}

\begin{thm}\label{thm: infinite sequences}
    Let $u$ be a strongly irreducible sequence on $r$ letters that contains each of its letters more than once. Then there exists a doubly infinite sequence on the alphabet $\{a_1,\dots,a_{r-1}\}\cup \Z$ that is $u$-saturated.
\end{thm}

\begin{proof}
    The idea is similar to the proof of \cref{thm: aa...bb}. Define $t_r(k)$ to be the infinite sequence over the alphabet $\{a_1,\dots, a_{r-1},x\}$ given by
    \[
    \dots (a_1\dots a_{r-1})(a_1\dots a_{r-1})(a_1\dots a_{r-1}x)^k (a_1\dots a_{r-1})(a_1\dots a_{r-1})\dots .
    \]
    Since $u$ contains each of its letters at least twice, the sequence $t_r(1)$ avoids $u$. Moreover, for sufficiently large $k$, it is clear that $t_r(k)$ contains a copy of $u$. Let $p$ be the maximum integer such that $t_r(p)$ avoids $u$. Note that $p\geq 1$.  
    
    For $n\in \mathbb Z$, define
    \[
    t[n] = (a_1\dots a_{r-1}n)^p.
    \]
    We then set
    \[
    s = \dots t[-1]t[0]t[1]t[2]\ldots = \bigcup_{n=-\infty}^\infty t[n].
    \]
    We claim that $s$ is $u$-saturated. It is clearly $r$-sparse.
    
    \smallskip \noindent\textit{Avoidance.} Suppose $s$ contains a copy of $u$. Since $u$ is strongly irreducible, such a copy can involve at most one letter from $\mathbb Z$. Thus it must lie on the alphabet $\{a_1,\dots, a_{r-1},n\}$ for some $n\in\mathbb Z$. But the subsequence of $s$ on these letters is isomorphic to $t_r(p)$, which by definition avoids $u$. Hence $s$ avoids $u$.
    
    \smallskip \noindent\textit{Saturation.} We cannot insert any $a_i$ without breaking $r$-sparsity. If we insert an integer $n$ without violating $r$-sparsity, then $s$ acquires a subsequence isomorphic to $t_r(p+1)$ over the alphabet $\{a_1,\dots, a_{r-1},n\}$. Since $p$ was chosen maximally, $t_r(p+1)$ contains a copy of $u$. It follows that $s$ is $u$-saturated.
\end{proof}

\end{document}
